\section{Fixed geometric auxiliary classes at the ramified CM prime}
\label{sec:cubic-geometric-route-diagnostic}

This section diagnoses a specific attempted replacement for the
weight-one Euler specialization. It proves no new mixed Selmer
vanishing. The coefficients and duals below are exactly those of
(CER.2), rather than the central self-dual coefficients of a
diagonal-cycle theorem.

\begin{proposition}[AC-G69-local: the de Rham singular quotient]
\label{AC-G69-local}
Let \(F/\mathbb Q_p\) and \(E/\mathbb Q_p\) be finite, let
\(W\) be an \(E\)-linear de Rham representation of \(G_F\),
and put \(U=W^*(1)\). There is a natural perfect pairing
\[
 \frac{H_g^1(F,U)}{H_f^1(F,U)}\ \times
 \frac{D_{\rm cris,F}(W)}{(1-\varphi)D_{\rm cris,F}(W)}
       \longrightarrow E .                             \tag{CGD.1}
\]
In particular the first space is zero if
\(D_{\rm cris,F}(W)^{\varphi=1}=0\).

For the actual Type II point of AC-G68, put \(p=269\),
\(F=K_p\), and
\[
 \begin{aligned}
 R&=V_f|_{G_K}\otimes\chi^j,\qquad U_K=R^*(1),\\
 W&=\operatorname{Ind}_{G_K}^{G_{\mathbb Q}}R
     =V_f\otimes\operatorname{Ind}_{G_K}^{G_{\mathbb Q}}\chi^j,\\
 U_{\mathbb Q}&=W^*(1).
 \end{aligned}                            \tag{CGD.2}
\]
Then
\[
 \begin{split}
 D_{\rm cris,F}(R)&=D_{\rm cris,F}(U_K)=0,\\
 D_{\rm cris,\mathbb Q_p}(W)&=
          D_{\rm cris,\mathbb Q_p}(U_{\mathbb Q})=0.
 \end{split}                                             \tag{CGD.3}
\]
For either local coefficient \(U_K\) over \(F\), or
\(U_{\mathbb Q}\) over \(\mathbb Q_p\), one has
\[
 H_e^1=H_f^1=H_g^1,\qquad
       \dim_E H_f^1=2,\qquad
       \dim_E(H^1/H_f^1)=2.                              \tag{CGD.4}
\]
Consequently an actual coefficient map applied to a class
whose localization is known to be de Rham has zero image in
this singular quotient. Such fixed-level classes cannot
supply its two independent singular directions.
\end{proposition}
\begin{proof}
Bloch--Kato, \emph{\(L\)-functions and Tamagawa numbers of
motives}, Proposition~3.8, printed pp.~354--359, makes
\(H_g^1(F,U)\) and \(H_e^1(F,W)\) exact annihilators,
and likewise the two finite local conditions. Taking the
induced perfect pairing on the corresponding quotients and
using Corollary~3.8.4, printed pp.~355--356, gives (CGD.1).
The map \(1-\varphi\) is an \(E\)-linear endomorphism
of a finite-dimensional \(E\)-space, even though
\(\varphi\) is semilinear on the unramified base-field
factor. Its cokernel and kernel have the same dimension.
This proves the first assertion and shows why the criterion
is not particular to \(269\).

For (CGD.3), AC-G68-point supplies the potentially crystalline
local representation \(V_f\), with finite inertial type
\(\tau\oplus\tau^{-1}\), where
\(\tau|_{I_p}=\psi|_{I_p}\) has order dividing two.
The restriction of \(\chi\) to inertia at \(269\) has
order three. Thus the two finite inertia characters of
\(R\) are \(\tau\chi^j\) and \(\tau^{-1}\chi^j\),
and neither is trivial. The same is true after duality or
Tate twist. On the induced coefficient the four inertia
characters are obtained by using \(\chi^j\) and
\(\chi^{-j}\), and again none is trivial.

Choose a finite Galois local extension on which the
representations are crystalline. Inertia in its descent
group acts on the full crystalline period module by exactly
these finite inertial types, and fixes the maximal unramified
coefficient field. Its invariant subspace is zero.
Crystalline periods over the original field are a subspace
of those invariants. This proves all four vanishings, with
no inference from ordinary Hodge--Tate weights alone.

The same crystalline vanishings imply that the zeroth
cohomology of each coefficient and its Tate dual is zero:
a fixed vector would give a nonzero crystalline period.
The Hodge--Tate weights of \(U_K\), with the convention
that \(\kappa_p\) has weight one, are
\(\{1-e,1+e\}\) at each of the two embeddings of \(F\).
Thus \(\dim_E D_{\rm dR,F}(U_K)/\operatorname{Fil}^0=2\),
also when \(e=1\); in that case the zero weight belongs
to \(\operatorname{Fil}^0\). The same dimension for
\(U_{\mathbb Q}\) follows from the four induced weights.
Corollary~3.8.4 gives the finite dimension in (CGD.4).
The local Euler characteristic gives total dimension four.
Using (CGD.1) and using that same corollary for \(U\)
proves both equalities in (CGD.4).

Finally \(H_g^1\) is the kernel of the coefficient map to
\(H^1(B_{\rm dR}\otimes- )\), as in Bloch--Kato
(3.7.2), printed p.~353. Every actual Galois-equivariant
coefficient map commutes with that map, and therefore
preserves \(H_g^1\). Its image here lies in \(H_f^1\).
This last assertion requires the stated de Rham property of
the originating class. There is a precise published source
for that property for the relevant fixed regulator classes:
Nekov\'a\v{r}--Nizio\l, \emph{Syntomic cohomology and
\(p\)-adic regulators for varieties over \(p\)-adic fields},
Algebra \& Number Theory \textbf{10} (2016),
Theorem~B, printed p.~1704, and Section~5B,
printed pp.~1778--1779. For \(2r-i-1\ge0\), their
Soul\'e regulator from \(K_{2r-i-1}(X)_0\) has target
\(H^1(F,H^i_{\rm et}(X_{\overline F},\mathbb Q_p(r)))\)
and is de Rham. An actual coefficient projection of that
regulator therefore has zero singular image here.
This uses the published syntomic comparison theorem;
its application permits bad reduction and does not require
an integral model. KLZ Proposition~3.3.3, printed p.~15,
applies this theorem to its fixed Rankin--Eisenstein classes,
before imposing the good-reduction assumption in the next
paragraph.
\end{proof}

\begin{proposition}[AC-G69-parity: fixed auxiliary CM quotients]
\label{AC-G69-parity}
Let \(f\) have the actual odd weight \(k=2e+1\), and
let \(g,h\) be algebraic CM forms of weights
\(l,m\ge2\). If an actual rank-one quotient of
\((M_g\otimes M_h)|_{G_K}\) is a finite Artin character
followed by an integral Tate twist, then \(l+m\) is even.
For these weights the diagonal algebraic invariant
\[
 \left(\operatorname{Sym}^{k-2}E^2\otimes
       \operatorname{Sym}^{l-2}E^2\otimes
       \operatorname{Sym}^{m-2}E^2\right)^{\operatorname{SL}_2}
                                                               \tag{CGD.5}
\]
is zero, with any determinant twist. Hence the usual fixed
balanced diagonal cycle cannot be turned into the actual
finite cubic coefficient through this rank-one quotient.

More generally, suppose \(X\) is smooth projective over a
number field \(F\) containing \(K\), and \(c\) is an integer. There is no nonzero
Galois-equivariant coefficient map
\[
 H^{2c-1}_{\rm et}(X_{\overline F},E)(c)
                    \longrightarrow U_K|_{G_F},        \tag{CGD.6}
\]
where the target is the actual modular coefficient in
(CGD.2). This includes a fixed ordinary Abel--Jacobi class
followed by an actual pure coefficient quotient.
\end{proposition}
\begin{proof}
Use geometric Frobenius weights, so that \(E(n)\) has
weight \(-2n\). The cohomological modular motive
\(M_g\) has weight \(l-1\), and \(M_h\) has weight
\(m-1\). Thus their tensor and every nonzero pure
rank-one quotient have weight \(l+m-2\). An Artin--Tate
coefficient has even weight. This proves the first parity
assertion. The central element \(-I\) of
\(\operatorname{SL}_2\) acts on the space before taking
invariants in (CGD.5) by
\((-1)^{k+l+m-6}=-1\). It acts trivially on every
determinant twist. There can therefore be no invariant.
This is a direct algebraic obstruction, independent of
whether the prime is split or inert.

The normalization must also be checked before using an
Abel--Jacobi theorem. In (CER.2),
\(V_f=M_f(e)\), and \(M_f\) has weight \(2e\).
Therefore \(V_f\), \(R\) and \(W\) have weight zero,
while \(U_K=R^*(1)\) has weight \(-2\).
The source of (CGD.6) has weight
\((2c-1)-2c=-1\). Choose a common good finite place
away from \(p\). Its Frobenius eigenvalues on these two
representations have complex absolute values
\((\#k_v)^{-1/2}\) and \((\#k_v)^{-1}\), respectively.
They are disjoint. An intertwiner of two finite-dimensional
operators with disjoint spectra is zero: their minimal
polynomials are coprime, so Bezout's identity applied to
an intertwiner proves this immediately. Every Galois map
in (CGD.6) is such an intertwiner. This proves the claim.

This argument uses the usual purity of the cohomological
modular motive and of smooth projective cohomology:
Deligne, \emph{La conjecture de Weil I}, Publications
Math\'ematiques de l'IH\'ES \textbf{43} (1974),
Theorem~1.6, printed p.~276,
\url{https://www.numdam.org/item/PMIHES_1974__43__273_0.pdf}.
The smooth projective hypothesis in (CGD.6) is exactly the
hypothesis of that published theorem. Its complete upstream
proof is not newly reproved here. The exact
cohomological convention for \(M_f\) is KLZ,
arXiv:1503.02888v6, Definition~2.8.1 and Remark~2.8.2,
printed pp.~13--14. The purity argument alone does not rule out a higher-Chow
coefficient, whose weight can be \(-2\). If its fixed
class belongs to the Soul\'e regulator family just described,
the separate local argument still gives zero singular image.
Neither argument rules out a nongeometric interpolation
specialization or another Euler-system method. In particular
it does not
apply the central self-dual diagonal-cycle coefficient to
our different, weight \(-2\), target.
\end{proof}

\paragraph{Exact scope of the newly inspected inert-prime source.}
Marannino, \emph{Diagonal cycles and anticyclotomic twists of
modular forms at inert primes}, Research in the Mathematical
Sciences \textbf{13} (2026), article~70,
\url{https://doi.org/10.1007/s40687-026-00651-w},
uses an even-weight form with trivial nebentypus, and
Assumption~1.1 requires residual absolute irreducibility,
\(p\)-distinction, and conductor prime to \(p\).
Theorem~1.2 is a rank-one result. Its CM coefficient
patching in Theorem~3.5 assumes \(p>\nu\) and
\(p\nmid\mathfrak f\); Remark~3.6 only relaxes the
latter at a split ordinary prime. The parity condition
in Section~3.2 agrees with (CGD.5).
Thus that theorem is not a vanishing theorem for the
present odd-weight, residually reducible,
conductor-\(269\) coefficient. Its Sections~1.1--1.2
explicitly distinguish horizontal classes from the
missing inert vertical interpolation.

\paragraph{Remaining construction.}
The auxiliary-prime norm relation by itself does not
supply a weight-one coefficient quotient. KLZ (3.2.1)
and Definition~3.3.1 require \(k,k'\ge0\), whereas
our second index is \(k'=-1\). Universal
Rankin--Iwasawa classes precede this specialization,
but the interpolating Theorem~6.3.4 uses ordinary
projectors and inverse unit-root operators. The
actual \(U_{269}=0\) is established in AC-K29-269.
An alternative must supply actual classes for (CGD.2),
with the required auxiliary-ray norms and a fixed
lattice denominator, and two independent singular
local images. No such construction follows from the
fixed geometric quotients tested here. The
already proved good-prime weight-one specializations
are outside the original motivic range and are
unaffected by this diagnosis.
